\documentclass[12pt]{article}
\usepackage[utf8]{inputenc}
\usepackage[T1]{fontenc}
\usepackage[margin=1in]{geometry}
\usepackage{setspace}
\usepackage{xcolor}
\usepackage[colorlinks=true, urlcolor=blue!60!black, linkcolor=blue!60!black]{hyperref}
\setlength{\parindent}{2em}
\setlength{\parskip}{1.25em}
\renewcommand{\baselinestretch}{1.0}
\begin{document}
\begin{center}
{\LARGE \textit{Allen v.\ Cooper}, 589 U.S.\ 248 (2020)\par}
\vspace{0.8em}
{\large Case Brief --- Four-Page Edition\par}
\end{center}
\vspace{1em}

\section{Memory Jogger}

In a unanimous judgment arising from North Carolina's unauthorized online posting of shipwreck videos of Blackbeard's \textit{Queen Anne's Revenge}, the Court held that Congress lacked authority under either Article I's Intellectual Property Clause or Section~5 of the Fourteenth Amendment to abrogate state sovereign immunity in the Copyright Remedy Clarification Act of 1990 (\href{https://supreme.justia.com/cases/federal/us/589/18-877/}{Justia}).

\section{Detailed Case Facts}

In 1718, Blackbeard's flagship \textit{Queen Anne's Revenge} ran aground off Beaufort, North Carolina. Salvage company Intersal discovered the wreck in 1996; under federal and state law it belonged to North Carolina. Intersal hired videographer Frederick Allen to document the salvage, and over more than a decade he produced videos and photographs and registered federal copyrights in them. In 2013 the State settled Allen's complaint about unauthorized online uploads for \$15,000, but it then posted five more of his videos and published a photograph in a newsletter. In 2015 the General Assembly enacted ``Blackbeard's Law,'' declaring shipwreck photographs and videos in state custody ``public record'' materials. Allen sued Governor Cooper and state officials for copyright infringement and money damages under the CRCA, which provides that a State ``shall not be immune'' from such suits. Congress had modeled the Patent Remedy Act on the CRCA, and \textit{Florida Prepaid} (1999) invalidated that patent twin.

\section{Procedural History}

The Eastern District of North Carolina denied the State's motion to dismiss, holding that Congress validly abrogated immunity under Section~5 based on a ``pattern'' of abusive state infringement. On interlocutory appeal the Fourth Circuit reversed, holding \textit{Florida Prepaid} barred abrogation under Section~5 as well as Article~I. The Supreme Court granted certiorari and affirmed on March~23, 2020.

\section{Judicial Votes}

Affirmed, \textbf{9--0}.
\begin{itemize}
\item \textbf{Majority:} Justice Kagan, joined by Chief Justice Roberts and Justices Alito, Sotomayor, Gorsuch, and Kavanaugh, and by Justice Thomas except for two paragraphs (stare decisis ``special justification'' and advice on future legislation).
\item \textbf{Concurring in part and in the judgment:} Justice Thomas.
\item \textbf{Concurring in the judgment:} Justice Breyer, joined by Justice Ginsburg.
\end{itemize}

\section{Holding}

North Carolina is immune from Allen's federal copyright suit for money damages. Although the CRCA clearly expressed Congress's intent to abrogate, Article~I's Intellectual Property Clause cannot abrogate state sovereign immunity under \textit{Seminole Tribe} and \textit{Florida Prepaid}, and \textit{Katz} is limited to the Bankruptcy Clause. Under Section~5 and \textit{City of Boerne}, a blanket abrogation for all copyright infringement fails the ``congruence and proportionality'' test where the record shows only a dozen possible examples of state infringement, mostly negligent, and no pattern of intentional infringement lacking adequate state remedies.

\section{Analysis of Opinions}

\textbf{Majority (Kagan).} ``\textit{Florida Prepaid} all but prewrote our decision today.'' On Article~I, the Intellectual Property Clause covers patents and copyrights alike, and \textit{Katz} reflects bankruptcy's \textit{in rem} jurisdiction and unique Founding-era history, where the States surrendered immunity in the ``plan of the Convention.'' Overruling \textit{Florida Prepaid} would require a ``special justification'' Allen did not offer. On Section~5, a State violates due process only when infringement is intentional or reckless and no adequate state remedy exists. The Oman Report showed twelve possible instances, only two apparently intentional, and testimony that state infringement mostly involved ``honest mistakes'' or ``sloppiness.'' The CRCA was ``out of proportion'' to any documented violation, though Congress could enact a tailored statute on a better record.

\textbf{Thomas (concurring in part and in the judgment).} He rejected the ``special justification'' standard, would overrule any ``demonstrably erroneous'' precedent, declined to advise Congress, and flagged as open whether copyrights are Fourteenth Amendment ``property.''

\textbf{Breyer, with Ginsburg (concurring in the judgment).} The Court ``went astray'' in \textit{Seminole Tribe} and erred in \textit{Florida Prepaid}, but stare decisis controls; he hoped Congress would venture into the ``great constitutional unknown'' to enact a tailored remedy.

\section{Examples: Future Applications}

\begin{enumerate}
\item \textbf{Same Side: Course-Portal Copying.} A publisher sues a state university for CRCA damages after a professor uploads scanned chapters. The university is an arm of the State and negligent academic copying cannot trigger federal damages liability.
\item \textbf{Same Side: Trade-Secret Abrogation.} Congress abrogates immunity for trade-secret claims under the Commerce Clause with no record of intentional state deprivations. Article~I cannot support it, and the untailored statute fails \textit{Boerne}.
\item \textbf{Opposite Side: \textit{Ex parte Young}.} A photographer seeks prospective injunctive relief against a state tourism director to stop ongoing display of her photos. \textit{Allen} does not disturb such suits.
\item \textbf{Opposite Side: Municipal Defendant.} A videographer sues a city transit authority for damages. Sovereign immunity protects States and arms of the State, not cities.
\item \textbf{Fence-Sitter: A Willful Piracy Act.} Congress authorizes damages only for willful state infringement without an adequate state remedy, backed by a 2021 Copyright Office study. It follows Kagan's roadmap, but whether the record suffices, and whether Thomas's ``property'' reservation prevails, is close.
\end{enumerate}

\section{Critique}

\textbf{Doctrinal criticism.} Vikram David Amar argues that the Court measures the overall frequency of state infringement rather than the proportion of covered conduct that is unconstitutional, and faults the revived ``special justification'' requirement where States have no legitimate reliance interest in pirating copyrights. Tran and Baker note that losing damages guts the incentive to enforce, leaving costly \textit{Ex parte Young} suits, and the ``good-for-one-clause-only'' reading of \textit{Katz} collapsed two years later in \textit{Torres}.

\textbf{Progressive and originalist views.} Progressives fault \textit{Allen} for entrenching an ahistorical expansion of sovereign immunity and, by pairing federal-court exclusivity with immunity, creating a right without a damages remedy. Originalists divide: Justice Thomas's camp supports immunity as a structural principle, while others argue neither \textit{Seminole Tribe}'s nontextual bar nor \textit{Boerne}'s test tracks original meaning.

\section{Key Quotations}

\begin{quote}
``But everything in \textit{Katz} is about and limited to the Bankruptcy Clause; the opinion reflects what might be called bankruptcy exceptionalism.''
\end{quote}
\noindent --- Justice Kagan, Opinion of the Court.

\begin{quote}
``The slate on which we write today is anything but clean. \textit{Florida Prepaid}, along with other precedent, forecloses each of Allen's arguments.''
\end{quote}
\noindent --- Justice Kagan, Opinion of the Court.

\begin{quote}
``That kind of tailored statute can effectively stop States from behaving as copyright pirates. Even while respecting constitutional limits, it can bring digital Blackbeards to justice.''
\end{quote}
\noindent --- Justice Kagan, Opinion of the Court.

\end{document}
